% Nouveau EGA V (cote 25): archive cote 25, site page 50 (PDF page 51).
% typescript. Draft compared: none.
% AI-transcribed by Claude Opus 5.5 (claude-opus-5-5, adaptive thinking, effort high), Message Batches API, batch ega25b (msgbatch_01A7n5YuXubUPNK1zsepi8RG), 2026-11-13; not yet checked by eye.
\scanpage{25/50}
\mspage{46}

\xout{\uline{Proposition} 8.16 On suppose f:X P non ramifié en}

Remords. Dans 8.13. et 8.14. , il suffit \inserted{de supposer} que f:X\hand{$\to$} P soit non
ramifié en les points lisses de X; pour le ``il suffit'', il suffit même
qu'il soit \inserted{non} ramifié sur un ouvert X' de X dont le complémentaire soit
de codimension \hand{$\geqslant$} \unsure{2}+1.

\uline{Proposition} 8.16. Supposons f:X\hand{$\to$} P non ramifié sur un ouvert
complémentaire d'un ensemble de codimension \hand{$\geqslant 2$,} X géom. normal et
de profondeur\hand{$\geqslant$} 3 en ses points fermés, enfin X géom. intègre et propre
sur k. Alors l'ensemble des \hand{$\zeta \in$} $\hand{$\check{\ }$}$P tels que \xout{$\phi$ soit $X_{k(\ )}$ régulier}
Y\hand{$_\zeta$} soit géom. normal et géom intègre de dimension dim X - 1 (est
constructible et) a un complémentaire de codimension \hand{$\geqslant$} 2.

En effet, par 8.14. b) il en est ainsi pour la propriété
`` Y\hand{$_\zeta$} est géom normal de dimension dim X - 1 '' (la propriété dimen
sionnelle exprime que \hand{$\phi_\zeta$} est X$_{k(\hand{$\zeta$})}$-régulier). Or en vertu de 6.1.
\uline{tous\xout{es}} les Y\hand{$_\zeta$} sont géom connexes. Or si Y\hand{$_\zeta$} est géom normal, il est
géom intègre sss il est géom connexe, ce qui achève la démonstration.

\uline{Remarques} 8.17. a) \inserted[typed]{L'hypothèse de 8.16. implique dim X \hand{$\geqslant$} 3.} Il est possible que, dès que X est géom \xout{intègre}irréd
et dim f(X) \hand{$\geqslant$} 3 (sans hypothèse de normalité ni de non ramification )
l'ensemble des \hand{$\zeta$} tels que Y\hand{$_\zeta$} soit géom irréd ait un complémentaire
de codimension \hand{$\geqslant$}2. On peut montrer en tous cas qu'il ne \unsure{contient}pas
d'hyperplan, \hand{cf. plus bas.}

b) La conclusion de 8.16. est en défaut si on abandonne
l'hypothèse prof$_x$X \hand{$\geqslant$} 3 pour x fermé, par exemple elle est en défaut pour
une quadrique non singulière \inserted{X} dans P$^3$ ; \deleted{dont} \inserted{\unsure{les sections par}} les hyperplans tangents
\xout{forment} sont \deleted{tous} réductibles (en fait, formés de couples de droites
concourrantes), et forment par ailleurs une famille de dimension 2
\margin{50} donc de codimension 1 dans $\hand{$\check{\ }$}$P (en fait, la duale de la quadrique \deleted{est} \deleted{\unsure{une}}
